% ============================================================================
%  Problems: ฟังก์ชันกำลังสอง (Quadratic Functions)
%  Ordered from easy (basic) → intermediate.
%  Shared by exercise.tex and answer-key.tex
% ============================================================================

% --- Part 1: พื้นฐาน (Basic) ---
\examsection{ตอนที่ 1: พื้นฐาน — รูปแบบและการจัดรูป}

\problem{
  จงระบุว่า $f(x) = -3x^2 + 5x - 1$ เป็นพาราโบลาหงายหรือคว่ำ
  พร้อมบอกจุดตัดแกน $y$
}{
  $a = -3 < 0$ \quad $\rightarrow$ \quad พาราโบลา \textbf{คว่ำ} $\cap$

  จุดตัดแกน $y$ คือ $f(0) = -1$ \quad ดังนั้น $(0, -1)$
}

\problem{
  จงเขียน $f(x) = x^2 - 8x + 10$ ในรูปจุดยอด $a(x-h)^2 + k$
  และบอกจุดยอด
}{
  $f(x) = (x^2 - 8x) + 10$

  $= (x^2 - 8x + 16) - 16 + 10$

  $= (x-4)^2 - 6$

  จุดยอด $= (4, -6)$
}

\problem{
  จงเขียน $f(x) = 3x^2 + 12x - 5$ ในรูปจุดยอด
  และบอกจุดยอด
}{
  $f(x) = 3(x^2 + 4x) - 5$

  $= 3(x^2 + 4x + 4) - 12 - 5$

  $= 3(x+2)^2 - 17$

  จุดยอด $= (-2, -17)$
}

\problem{
  จงหาแกนสมมาตร จุดยอด และค่าต่ำสุดของ $f(x) = x^2 - 10x + 7$
}{
  $h = -\frac{-10}{2(1)} = 5$

  $k = f(5) = 25 - 50 + 7 = -18$

  แกนสมมาตร $x = 5$, จุดยอด $(5, -18)$, ค่าต่ำสุด $= -18$
}

% --- Part 2: การประยุกต์ (Intermediate) ---
\examsection{ตอนที่ 2: การประยุกต์ — ราก, ผลบวก/ผลคูณ, อสมการ}

\problem{
  จงหาจำนวนรากของสมการ $3x^2 - 4x + 2 = 0$ โดยใช้ discriminant
}{
  $\Delta = (-4)^2 - 4(3)(2) = 16 - 24 = -8$

  $\Delta < 0$ \quad ดังนั้น ไม่มีรากจริง
}

\problem{
  จงหารากของสมการ $2x^2 - 5x - 3 = 0$
}{
  $\Delta = 25 - 4(2)(-3) = 25 + 24 = 49$

  $x = \frac{5 \pm 7}{4}$ \quad $\rightarrow$ \quad $x = 3$ หรือ $x = -\frac{1}{2}$
}

\problem{
  จงหาผลบวกและผลคูณของรากของสมการ $4x^2 + 8x - 12 = 0$
  โดยไม่ต้องแก้สมการ
}{
  $r_1 + r_2 = -\frac{8}{4} = -2$

  $r_1 r_2 = \frac{-12}{4} = -3$
}

\problem{
  จงสร้างสมการกำลังสองที่มีรากเป็น $-2$ และ $6$ โดยมีสัมประสิทธิ์นำเป็น $1$
}{
  $r_1 + r_2 = -2 + 6 = 4$ \qquad
  $r_1 r_2 = (-2)(6) = -12$

  $x^2 - 4x - 12 = 0$
}

\problem{
  จงแก้อสมการ $x^2 - x - 6 < 0$ และเขียนคำตอบในรูปช่วง
}{
  $(x-3)(x+2) < 0$ \quad รากคือ $x = -2, 3$

  $a=1>0$ (หงาย) — ต่ำกว่า $0$ ระหว่างราก

  $x \in (-2, 3)$
}

\problem{
  ร้านขายขนมมีกำไรต่อวัน $P(x) = -5x^2 + 100x - 200$ บาท
  เมื่อขายได้ $x$ กล่อง จงหาจำนวนกล่องที่ทำให้ได้กำไรสูงสุด
  และกำไรสูงสุดเป็นเท่าใด
}{
  $a = -5 < 0$ (คว่ำ มีค่าสูงสุด)

  $x = -\frac{100}{2(-5)} = 10$ กล่อง

  $P(10) = -5(100) + 100(10) - 200 = -500 + 1000 - 200 = 300$ บาท

  ขาย $10$ กล่อง ได้กำไรสูงสุด $300$ บาท
}

% --- Part 3: ระดับ สอวน. (POSN Level) ---
\examsection{ตอนที่ 3: ระดับข้อสอบ สอวน.}

\problem{
  (ข้อสอบจริงปี 68 ข้อ 27)  โรงงานแห่งหนึ่งรับจ้างผลิตสินค้า ซึ่งในการผลิตสินค้าจำนวน $x$ ชิ้นนั้น จะมีต้นทุนในการผลิตทั้งหมดเท่ากับ $550,000 + 200x$ บาท โดยจะจำหน่ายสินค้าในราคาชิ้นละ $1,800 - x$ บาท ถ้าสินค้าสามารถขายได้ทั้งหมด แล้วจะทำกำไรสูงสุดเมื่อขายสินค้าได้กี่ชิ้น กำไรสูงสุดกี่บาท จะมีกำไรเมื่อผลิตสินค้าในช่วงใด
}{
  กำไรสูงสุดเมื่อผลิตและขาย 800 ชิ้น \\
  กำไรสูงสุด = 90,000 บาท \\
  มีกำไร เมื่อผลิตสินค้าในช่วง 500 < x < 1100
}

\problem{
  (ข้อสอบจริงปี 67 ข้อ 14) ถ้ากราฟพาราโบลา $y=ax^2+bx+c$ มีจุดสูงสุดของกราฟ คือ $(-2,-3)$ และมีจุดตัดแกน $y$ ที่ $(0,-7)$ จงหาค่า $y$ เมื่อ $x=-5$
}{
  $-12$
}

\problem{
  (ข้อสอบจริงปี 67 ข้อ 15)  ต้นทุนการผลิตเสื้อจำนวน $q$ ตัว เป็นเงิน $q^2-40q+400$ บาท ถ้าขายเสื้อตัวละ 60 บาท แล้ว จำนวนเสื้อที่จะผลิตเพื่อให้ได้กำไรสูงสุด และกำไรรวมเท่ากับเท่าไร
}{
  เสื้อ 50 ตัว กำไร 2,100 บาท
}

\problem{
  (ข้อสอบจริงปี 66 ข้อ 19) กำหนด พาราโบลาหงายรูปหนึ่งมีจุดตัดแกน $x$ ที่ $(1,0)$ และ $(3,0)$ และตัดกับเส้นตรง $x+2y=-2$ ที่ $x=2$ จงหาจุดตัดแกน $y$ ของพาราโบลารูปนี้ 
}{
  $(0,6)$
}